\documentclass[11pt, a4paper, oneside]{article}
\usepackage{xeCJK}
\usepackage{amsmath, amsthm, amssymb, bm, color, framed, graphicx, hyperref, mathrsfs}
\usepackage{geometry}
\usepackage{fancyhdr}
\usepackage{enumitem}
\usepackage{lastpage}
\usepackage{booktabs}
\usepackage[raggedright]{titlesec}
\titleformat{\section}
  {\normalfont\bfseries}
  {\thesection}
  {1em}
  {}

\geometry{a4paper, left=2.5cm, right=2.5cm, top=2.5cm, bottom=2.5cm}
\setlength{\headheight}{14pt}

\setlength{\parindent}{0pt}

\fancypagestyle{plain}{
    \fancyhf{}
    \fancyhead[L]{Assignment 4}
    \fancyhead[R]{CS224n Natural Language Processing}
    \fancyfoot[C]{\thepage}
}

\pagestyle{plain}

\title{\textbf{Assignment 4}}
\author{}
\date{\empty}

\begin{document}

\maketitle
\section{Attention Exploration (14 points)}
\begin{enumerate}[label=(\alph*), left=0pt, labelsep=1em, itemsep=1em]
    \item 
    \begin{enumerate}[label=\roman*., left=0pt, labelsep=1em, itemsep=1em]
        \item The categorical distribution $\alpha$ puts almost all of its
        weight on $\alpha_j$ when the score $k_j^\top q$ is much larger than
        every other score $k_i^\top q$; equivalently,
        $\sum_{i\neq j}\exp(k_i^\top q-k_j^\top q)\ll 1$.
        \item Under this condition, $\alpha_j\approx 1$ and
        $\alpha_i\approx 0$ for every $i\neq j$. Therefore,
        $c=\sum_{i=1}^n\alpha_i v_i\approx v_j$.
    \end{enumerate}
    \item Let $\ell>0$ be large and choose
    \[
        q=\ell(k_a+k_b).
    \]
    Since the keys are orthogonal unit vectors,
    $k_a^\top q=k_b^\top q=\ell$, whereas $k_i^\top q=0$ for
    $i\notin\{a,b\}$. Hence
    \[
        \alpha_a=\alpha_b
        =\frac{e^\ell}{2e^\ell+(n-2)}
        =\frac{1}{2+(n-2)e^{-\ell}}\approx\frac12
    \]
    whenever $e^\ell\gg n-2$. The remaining attention weights are negligible,
    so $c\approx\frac12(v_a+v_b)$.
    \item \label{1c}
    \begin{enumerate}[label=\roman*., left=0pt, labelsep=1em, itemsep=1em]
        \item To distinguish the covariance scale from the attention weights
        $\alpha_i$, let $\varepsilon$ denote the scalar called $\alpha$ in the
        question. Thus, $\Sigma_i=\varepsilon I$ with $\varepsilon\to0$.
        Choose
        \[
            q=\ell(\mu_a+\mu_b),
            \qquad e^\ell\gg n-2.
        \]
        Write $k_i=\mu_i+\delta_i$, where
        $\delta_i\sim\mathcal N(0,\varepsilon I)$. Since the means are
        orthogonal unit vectors,
        \[
            k_a^\top q\approx k_b^\top q\approx\ell,
            \qquad
            k_i^\top q\approx0\quad(i\notin\{a,b\}).
        \]
        Therefore, $\alpha_a\approx\alpha_b\approx\frac12$ and
        $c\approx\frac12(v_a+v_b)$. To make the conditions on $\ell$ precise,
        first note that the random score perturbation $\delta_i^\top q$ is a
        linear projection of a Gaussian vector and is therefore Gaussian. Its
        mean and variance are
        \[
            \mathbb E[\delta_i^\top q]=0,
            \qquad
            \operatorname{Var}(\delta_i^\top q)
            =q^\top(\varepsilon I)q
            =\varepsilon\lVert q\rVert^2
            =2\varepsilon\ell^2,
        \]
        where $\lVert q\rVert^2=\ell^2\lVert\mu_a+\mu_b\rVert^2
        =2\ell^2$. Hence
        \[
            \delta_i^\top q\sim
            \mathcal N(0,2\varepsilon\ell^2),
        \]
        whose standard deviation is $\sqrt{2}\ell\sqrt{\varepsilon}$.
        We may thus choose
        \[
            \log(n-2)\ll\ell\ll\varepsilon^{-1/2}.
        \]
        The lower bound implies $e^\ell\gg n-2$, suppressing the total
        contribution of the irrelevant keys. The upper bound is equivalent to
        $\ell\sqrt{\varepsilon}\ll1$, ensuring that the random score
        perturbations are small and that the scores of $k_a$ and $k_b$ remain
        approximately equal.

        \item Now
        $\Sigma_a=\varepsilon I+\frac12\mu_a\mu_a^\top$. Ignoring the
        vanishing orthogonal noise, write $k_a\approx r\mu_a$, where the random
        magnitude
        \[
            r=\mu_a^\top k_a
            \sim\mathcal N\!\left(1,\varepsilon+\frac12\right).
        \]
        Using the query from part (i), the two relevant scores satisfy
        \[
            q^\top k_a\approx\ell r,
            \qquad q^\top k_b\approx\ell,
            \qquad
            \frac{\alpha_a}{\alpha_b}\approx e^{\ell(r-1)}.
        \]
        Thus, the softmax exponentially amplifies fluctuations in $r$. When
        $r>1$, $c$ is close to $v_a$; when $r<1$, $c$ is close to $v_b$; only
        when $r$ is very close to $1$ is $c$ close to
        $\frac12(v_a+v_b)$. Consequently, individual samples produce outputs
        that vary substantially between vectors near $v_a$ and $v_b$, so the
        variance of $c$ is much larger than in part (i), even though its average
        over the symmetric fluctuations may remain near
        $\frac12(v_a+v_b)$.
    \end{enumerate}

    \item
    \begin{enumerate}[label=\roman*., left=0pt, labelsep=1em, itemsep=1em]
        \item Let $\ell>0$ be large and choose two different queries
        \[
            q_1=\ell\mu_a,
            \qquad
            q_2=\ell\mu_b.
        \]
        When $\Sigma_i=\varepsilon I$ and $\varepsilon\to0$, we have
        $k_i\approx\mu_i$. Thus, for $q_1$, the score of $k_a$ is
        approximately $\ell$, while all other scores are approximately zero;
        similarly, $q_2$ selects $k_b$. If $e^\ell\gg n-1$ (and
        $\ell\sqrt{\varepsilon}\ll1$), then
        \[
            c_1\approx v_a,
            \qquad
            c_2\approx v_b,
            \qquad
            c=\frac12(c_1+c_2)\approx\frac12(v_a+v_b).
        \]

        \item With
        $\Sigma_a=\varepsilon I+\frac12\mu_a\mu_a^\top$, write
        $k_a\approx r\mu_a$, where
        $r\sim\mathcal N(1,\varepsilon+\frac12)$. The scores relevant to the
        first head are
        \[
            q_1^\top k_a\approx\ell r,
            \qquad
            q_1^\top k_i\approx0\quad(i\neq a).
        \]
        Ignoring the cases $q_1^\top k_a<0$ as permitted, the first head still
        selects $v_a$ whenever $\ell r\gg\log(n-1)$. Hence
        $c_1\approx v_a$, although it has some residual variance when $r$ is
        positive but close to zero. For the second head,
        \[
            q_2^\top k_b\approx\ell,
            \qquad
            q_2^\top k_a\approx0,
        \]
        because $\mu_a^\top\mu_b=0$. Therefore, $c_2\approx v_b$ and has
        vanishingly small variance. Consequently,
        \[
            c=\frac12(c_1+c_2)\approx\frac12(v_a+v_b)
        \]
        across samples. In particular, since $c_2$ is nearly deterministic,
        \[
            \operatorname{Var}(c)
            \approx\frac14\operatorname{Var}(c_1),
        \]
        so the output is substantially more stable than the single-headed
        construction in part (c).
    \end{enumerate}

    \item Multi-headed attention assigns $v_a$ and $v_b$ to separate heads,
    so a change in the magnitude of $k_a$ affects mainly the head responsible
    for $v_a$ rather than changing the competition between $v_a$ and $v_b$
    within one softmax. Averaging the head outputs therefore preserves both
    values more reliably and reduces the variance of the final output.
\end{enumerate}
\newpage
\section{Position Embeddings Exploration (6 points)}
\begin{enumerate}[label=(\alph*), left=0pt, labelsep=1em, itemsep=1em]
    \item
    \begin{enumerate}[label=\roman*., left=0pt, labelsep=1em, itemsep=1em]
        \item Since
        $X_{\mathrm{perm}}=PX$, the query, key, and value matrices become
        \begin{align*}
            Q_{\mathrm{perm}}
            &=X_{\mathrm{perm}}W_Q=PXW_Q=PQ,\\
            K_{\mathrm{perm}}&=PK,\\
            V_{\mathrm{perm}}&=PV.
        \end{align*}
        Define $A=QK^\top/\sqrt d$. Then
        \begin{align*}
            H_{\mathrm{perm}}
            &=\operatorname{softmax}\!\left(
                \frac{Q_{\mathrm{perm}}K_{\mathrm{perm}}^\top}{\sqrt d}
              \right)V_{\mathrm{perm}}\\
            &=\operatorname{softmax}(PAP^\top)PV\\
            &=P\operatorname{softmax}(A)P^\top PV\\
            &=P\operatorname{softmax}(A)V
             =PH,
        \end{align*}
        where we used $P^\top P=I_T$. A permutation leaves the all-ones
        vector unchanged, so $P\mathbf1=\mathbf1$. Consequently,
        \begin{align*}
            Z_{\mathrm{perm}}
            &=\operatorname{ReLU}
              (H_{\mathrm{perm}}W_1+\mathbf1b_1)W_2+\mathbf1b_2\\
            &=\operatorname{ReLU}
              \bigl(PHW_1+P\mathbf1b_1\bigr)W_2+P\mathbf1b_2\\
            &=\operatorname{ReLU}
              \bigl(P(HW_1+\mathbf1b_1)\bigr)W_2+P\mathbf1b_2\\
            &=P\operatorname{ReLU}
              (HW_1+\mathbf1b_1)W_2+P\mathbf1b_2\\
            &=P\left[
                \operatorname{ReLU}(HW_1+\mathbf1b_1)W_2+\mathbf1b_2
              \right]
             =PZ.
        \end{align*}
        Therefore, permuting the input rows by $P$ permutes the output rows by
        the same matrix $P$.

        \item Although permutation equivariance is useful for unordered set
        inputs, it is problematic for text because word order carries
        syntactic and semantic information. Without positional information,
        permuting the input tokens merely permutes the output representations
        in the same way; after matching each output with its corresponding
        token, the representations are otherwise unchanged. Therefore, the
        Transformer cannot model order-dependent relationships, such as which
        noun is the subject and which is the object.
    \end{enumerate}

    \item
    \begin{enumerate}[label=\roman*., left=0pt, labelsep=1em, itemsep=1em]
        \item Yes. After permuting the tokens, the position-encoded input is
        \[
            X_{\mathrm{pos,perm}}=PX+\Phi,
        \]
        because $\Phi$ remains attached to the fixed sequence positions. In
        general,
        \[
            PX+\Phi\neq P(X+\Phi)=PX+P\Phi.
        \]
        Thus, the permutation-equivariance result from part (a) no longer
        applies. A token moved to a new position receives a different position
        embedding, allowing the Transformer to model word order.

        \item No. Suppose two integer positions $t$ and $t'$ had identical
        position embeddings. Taking $i=0$ gives
        \[
            \sin t=\sin t',
            \qquad
            \cos t=\cos t'.
        \]
        Hence $t-t'=2\pi m$ for some integer $m$. Since $t-t'$ is an integer,
        this is possible only when $m=0$, so $t=t'$. Therefore, two distinct
        positions cannot have identical position embeddings.
    \end{enumerate}
\end{enumerate}
\section{Pretrained Transformer models and knowledge access (35 points)}
\begin{enumerate}[label=(\alph*), left=0pt, labelsep=1em, itemsep=1em]
    \item[d] Finetune w/o pretraining: Correct: 11.0 out of 500.0: 2.1999999999999997\%.
    London baseline: Correct: 25.0 out of 500.0: 5.0\%.
    \item[f] Finetune w/ pretraining: Correct: 128.0 out of 500.0: 25.6\%
    \item[g]
    \begin{enumerate}[label=\roman*., left=0pt, labelsep=1em, itemsep=1em]
        \item
        For $k=1,\ldots,d/2$, represent the $k$-th pair of features as the
        complex number
        \[
            z_{t,k}=x_t^{(2k-1)}+\mathrm{i}x_t^{(2k)},
            \qquad
            \theta_k=10000^{-2(k-1)/d}.
        \]
        The $k$-th element of the element-wise product in Equation 4 is
        \[
            u_{t,k}
            =\bigl(\cos(t\theta_k)+\mathrm{i}\sin(t\theta_k)\bigr)
             \bigl(x_t^{(2k-1)}+\mathrm{i}x_t^{(2k)}\bigr).
        \]
        Expanding this product and using $\mathrm{i}^2=-1$ gives
        \begin{align*}
            u_{t,k}
            &=x_t^{(2k-1)}\cos(t\theta_k)
              +\mathrm{i}x_t^{(2k-1)}\sin(t\theta_k)\\
            &\quad+\mathrm{i}x_t^{(2k)}\cos(t\theta_k)
              -x_t^{(2k)}\sin(t\theta_k)\\
            &=\bigl[x_t^{(2k-1)}\cos(t\theta_k)
              -x_t^{(2k)}\sin(t\theta_k)\bigr]\\
            &\quad+\mathrm{i}\bigl[x_t^{(2k-1)}\sin(t\theta_k)
              +x_t^{(2k)}\cos(t\theta_k)\bigr].
        \end{align*}
        Therefore, the real and imaginary parts of $u_{t,k}$ satisfy
        \[
            \begin{bmatrix}
                \operatorname{Re}(u_{t,k})\\
                \operatorname{Im}(u_{t,k})
            \end{bmatrix}
            =
            \begin{bmatrix}
                \cos(t\theta_k)&-\sin(t\theta_k)\\
                \sin(t\theta_k)&\cos(t\theta_k)
            \end{bmatrix}
            \begin{bmatrix}
                x_t^{(2k-1)}\\
                x_t^{(2k)}
            \end{bmatrix}.
        \]
        This is exactly the $k$-th $2\times2$ rotation block in Equation 3.
        Stacking the real and imaginary parts for all $k$ in their original
        order gives
        \[
            \operatorname{RoPE}(x_t,t)
            =
            \begin{bmatrix}
                \operatorname{Re}(u_{t,1})\\
                \operatorname{Im}(u_{t,1})\\
                \vdots\\
                \operatorname{Re}(u_{t,d/2})\\
                \operatorname{Im}(u_{t,d/2})
            \end{bmatrix},
        \]
        which is the vector produced by the block-diagonal matrix in
        Equation 3. Thus, Equation 4 computes all two-dimensional rotations
        simultaneously by element-wise complex multiplication; in code, this
        requires reshaping the real vector into pairs, converting the pairs to
        complex numbers, applying the multiplication, converting back to real
        numbers, and flattening the pairs.

        \item
        For a two-dimensional vector represented by a complex number, the
        RoPE operation is
        \[
            \operatorname{RoPE}(z,t)=e^{\mathrm{i}t\theta}z,
        \]
        and the corresponding real dot product is
        \[
            \langle z_1,z_2\rangle
            =\operatorname{Re}\!\left(z_1\overline{z_2}\right),
        \]
        where the bar denotes complex conjugation. Therefore,
        \begin{align*}
            \left\langle \operatorname{RoPE}(z_1,t_1),
            \operatorname{RoPE}(z_2,t_2)\right\rangle
            &=\operatorname{Re}\!\left[
                \left(e^{\mathrm{i}t_1\theta}z_1\right)
                \overline{\left(e^{\mathrm{i}t_2\theta}z_2\right)}
              \right]\\
            &=\operatorname{Re}\!\left[
                e^{\mathrm{i}t_1\theta}z_1
                e^{-\mathrm{i}t_2\theta}\overline{z_2}
              \right]\\
            &=\operatorname{Re}\!\left[
                e^{\mathrm{i}(t_1-t_2)\theta}z_1\overline{z_2}
              \right].
        \end{align*}
        On the other hand, since
        $\operatorname{RoPE}(z_2,0)=e^{0}z_2=z_2$, we have
        \begin{align*}
            \left\langle \operatorname{RoPE}(z_1,t_1-t_2),
            \operatorname{RoPE}(z_2,0)\right\rangle
            &=\operatorname{Re}\!\left[
                \left(e^{\mathrm{i}(t_1-t_2)\theta}z_1\right)
                \overline{z_2}
              \right]\\
            &=\operatorname{Re}\!\left[
                e^{\mathrm{i}(t_1-t_2)\theta}z_1\overline{z_2}
              \right].
        \end{align*}
        The two expressions are identical. Hence,
        \[
            \left\langle \operatorname{RoPE}(z_1,t_1),
            \operatorname{RoPE}(z_2,t_2)\right\rangle
            =
            \left\langle \operatorname{RoPE}(z_1,t_1-t_2),
            \operatorname{RoPE}(z_2,0)\right\rangle,
        \]
        showing that the dot product depends on the positions only through
        their relative displacement $t_1-t_2$.
    \end{enumerate}
\end{enumerate}
\end{document}
